\clearpage
\section{Chapter 5: Thermodynamics \& Primordial Soup}

\subsection{Questions}

\begin{enumerate}[label=5.\arabic*]
  \item \label{q:5.1} \textbf{Thermal Distributions}

  Assume $\mu=0$ and $E=5T$.
  \begin{enumerate}[label=(\alph*)]
    \item Compute $f_{FD}=1/(e^{E/T}+1)$.
    \item Compute $f_{BE}=1/(e^{E/T}-1)$.
    \item Compare both with Maxwell-Boltzmann $f_{MB}=e^{-E/T}$ and explain why they are close in this limit.
  \end{enumerate}

  \item \label{q:5.2} \textbf{Relativistic Densities and Equation of State}

  For one relativistic boson species with degeneracy $g=2$:
  \begin{enumerate}[label=(\alph*)]
    \item Write $\rho(T)$ and $P(T)$.
    \item Evaluate $\rho$ at $T=1\,\mathrm{MeV}$ in natural units (leave answer in $\mathrm{MeV}^4$).
    \item Show that $w\equiv P/\rho=1/3$.
  \end{enumerate}

  \item \label{q:5.3} \textbf{Neutrino-Photon Temperature Split}

  \begin{enumerate}[label=(\alph*)]
    \item \textit{(Use the instantaneous neutrino decoupling approximation)} Using entropy conservation around $e^+e^-$ annihilation, derive
    \[
      \frac{T_\nu}{T_\gamma}=\left(\frac{4}{11}\right)^{1/3}.
    \]
    \item Compute the numerical value of this ratio.
    \item Using
    \[
      g_*=2+\frac{7}{8}\cdot 6\left(\frac{T_\nu}{T_\gamma}\right)^4,
      \quad
      g_{*s}=2+\frac{7}{8}\cdot 6\left(\frac{T_\nu}{T_\gamma}\right)^3,
    \]
    evaluate $g_*$ and $g_{*s}$ today.
  \end{enumerate}

  \item \label{q:5.4} \textbf{Modified Neutrino Decoupling: Massive Photon Scenario}

  Reconsider the neutrino-photon decoupling process (Question 5.3), but now examine a modified scenario where the photon has acquired a small mass (a pedagogical thought experiment). Before $e^+e^-$ annihilation, both neutrinos and the ``massive'' photon participate in the thermal bath alongside $e^\pm$ pairs. \textit{Note: Consider three families of neutrinos, as is standard in particle physics.}
  
  \begin{enumerate}[label=(\alph*)]
    \item 
    In the standard scenario (massless photon), count the degrees of freedom: photons ($g=2$), three neutrino families ($g_\nu=2$ each, so $3 \times 2 = 6$ total), and $e^\pm$ pairs ($g_{e^\pm}=4$). With the $\frac{7}{8}$ entropy weighting for fermions:
    \[
      g_{*s}^{\mathrm{before}} = 2 + \frac{7}{8}\times(6+4) = 2 + \frac{7}{8}\times 10 = 2 + 8.75 = 10.75.
    \]
    Now consider a ``massive photon'' scenario where the photon, as a massive vector boson, has $g_\gamma = 3$ degrees of freedom (one longitudinal + two transverse modes). Modify the calculation to find $g_{*s,1}^{\mathrm{before}}$ (massive photon case) and $g_{*s,2}^{\mathrm{after}}$ (after annihilation, when only the photon remains). Then derive the new neutrino-photon temperature ratio $T_\nu/T_\gamma$ under these conditions.
    
    \item If the baryon-to-photon ratio today is $\eta \equiv n_B/n_\gamma = 2.0 \times 10^{-10}$ (where $n_\gamma \propto T^3$), estimate the present-day CMB temperature $T_{\gamma}^{(0)}$ in this modified scenario and compare with the standard value of $T_{\gamma}^{(0)} \approx 2.725\,\mathrm{K}$. (Hint: Use $n_\gamma \approx 0.2437\, T_\gamma^3$ in natural units.)
  \end{enumerate}
\end{enumerate}

\bigskip
\bigskip
\bigskip

\subsection{Solutions}

\subsubsection*{Solution 5.1: Thermal Distributions}

\begin{enumerate}[label=(\alph*)]
  \item
  \[
    f_{FD}=\frac{1}{e^5+1}=\frac{1}{148.413+1}\approx 6.69\times 10^{-3}.
  \]

  \item
  \[
    f_{BE}=\frac{1}{e^5-1}=\frac{1}{148.413-1}\approx 6.78\times 10^{-3}.
  \]

  \item
  \[
    f_{MB}=e^{-5}\approx 6.74\times 10^{-3}.
  \]
  Since $E/T=5\gg1$, the $\pm1$ in FD/BE denominators is negligible, so both approach MB.
\end{enumerate}

\subsubsection*{Solution 5.2: Relativistic Densities and Equation of State}

\begin{enumerate}[label=(\alph*)]
  \item For relativistic bosons:
  \[
    \rho = \frac{\pi^2}{30}gT^4,
    \qquad
    P = \frac{1}{3}\rho.
  \]

  \item With $g=2$, $T=1\,\mathrm{MeV}$:
  \[
    \rho = \frac{\pi^2}{30}\cdot 2\cdot (1\,\mathrm{MeV})^4
    = \frac{\pi^2}{15}\,\mathrm{MeV}^4
    \approx 0.658\,\mathrm{MeV}^4.
  \]

  \item
  \[
    w\equiv\frac{P}{\rho}=\frac{1}{3}.
  \]
\end{enumerate}

\subsubsection*{Solution 5.3: Neutrino-Photon Temperature Split}

\begin{enumerate}[label=(\alph*)]
  \item Under the instantaneous neutrino decoupling approximation, entropy in the coupled $(\gamma,e^\pm)$ plasma is conserved:
  \[
    g_{*s}^{\mathrm{before}}T_{\mathrm{before}}^3a^3
    = g_{*s}^{\gamma,\mathrm{after}}T_\gamma^3a^3.
  \]
  Before annihilation, for the coupled plasma,
  \[
    g_{*s}^{\mathrm{before}} = 2 + \frac{7}{8}\times 4 = \frac{11}{2},
  \]
  after annihilation only photons remain coupled: $g_{*s}^{\gamma,\mathrm{after}}=2$.
  Since neutrinos are decoupled, $T_\nu\propto a^{-1}$ and effectively track $T_{\mathrm{before}}a$:
  \[
    \frac{T_\nu}{T_\gamma} = \left(\frac{2}{11/2}\right)^{1/3} = \left(\frac{4}{11}\right)^{1/3}.
  \]

  \item
  \[
    \left(\frac{4}{11}\right)^{1/3} \approx 0.714.
  \]

  \item
  \[
    g_* = 2 + \frac{7}{8}\cdot 6\left(\frac{4}{11}\right)^{4/3} \approx 3.36,
  \]
  \[
    g_{*s} = 2 + \frac{7}{8}\cdot 6\left(\frac{4}{11}\right) \approx 3.91.
  \]
  So today $g_*\neq g_{*s}$ because $T_\nu\neq T_\gamma$.
  \textit{Rigorous note:} beyond the instantaneous-decoupling approximation, small corrections from non-instantaneous decoupling and QED effects exist (often encoded via $N_\mathrm{eff}\simeq 3.046$), but they do not change the leading-order conclusion.
\end{enumerate}

\subsubsection*{Solution 5.4: Modified Neutrino Decoupling: Massive Photon Scenario}

\begin{enumerate}[label=(\alph*)]
  \item \textbf{Standard scenario (massless photon) with three neutrino families:}
  
  Count the relativistic degrees of freedom before $e^+e^-$ annihilation:
  \begin{itemize}
    \item Photons: $g_\gamma=2$
    \item Three neutrino families: each with $g_\nu=2$, so $3 \times 2 = 6$ total
    \item Electron-positron pairs: $g_{e^\pm}=4$
  \end{itemize}
  
  For entropy, fermions contribute with a factor $\frac{7}{8}$ relative to bosons:
  \[
    g_{*s}^{\mathrm{before,std}} = 2 + \frac{7}{8}\times(6+4) = 2 + \frac{7}{8}\times 10 = 2 + 8.75 = 10.75.
  \]
  
  After $e^+e^-$ annihilation, only photons remain (massless case):
  \[
    g_{*s}^{\mathrm{after,std}} = 2.
  \]
  
  \textbf{Massive photon scenario:}
  
  A massive photon (as a massive vector boson) has three polarization states instead of two: one longitudinal plus two transverse modes, so $g_\gamma^{\mathrm{massive}} = 3$.
  
  Before annihilation:
  \[
    g_{*s,1}^{\mathrm{before}} = 3 + \frac{7}{8}\times(6+4) = 3 + 8.75 = 11.75.
  \]
  
  After annihilation (only the massive photon):
  \[
    g_{*s,2}^{\mathrm{after}} = 3.
  \]
  
  From entropy conservation, the neutrino-photon temperature ratio is:
  \[
    \frac{T_\nu}{T_\gamma} = \left(\frac{g_{*s,2}^{\mathrm{after}}}{g_{*s,1}^{\mathrm{before}}}\right)^{1/3} = \left(\frac{3}{11.75}\right)^{1/3} = \left(0.2553\right)^{1/3} \approx 0.635.
  \]
  
  For comparison, the standard result with two families (as used in earlier solutions) would give $(2/10.75)^{1/3} \approx 0.613$, and the three-family standard massless case gives $(2/10.75)^{1/3} \approx 0.613$ as well. The massive photon scenario yields a slightly higher temperature ratio due to the additional photon polarization state after annihilation.

  \item The photon number density in thermal equilibrium scales as:
  \[
    n_\gamma(T) = 0.2437\, T_\gamma^3 \quad \text{(in natural units)}.
  \]
  
  For baryon-to-photon ratio $\eta = n_B/n_\gamma = 2.0 \times 10^{-10}$, this ratio is fixed by Big Bang Nucleosynthesis (BBN) observations and remains independent of the photon mass scenario in this thought experiment.
  
  The present-day CMB temperature is an empirically measured quantity from blackbody observations:
  \[
    T_\gamma^{(0)} \approx 2.725\,\mathrm{K}.
  \]
  
  In this pedagogical massive-photon scenario, the CMB temperature would remain essentially unchanged at $T_{\gamma}^{(0)} \approx 2.725\,\mathrm{K}$ because the temperature is determined by direct observation of the CMB blackbody spectrum, not derived from thermodynamic calculations alone.
  
  The key educational insight is:
  \begin{itemize}
    \item The **temperature ratio** $T_\nu/T_\gamma$ depends on the degrees of freedom at decoupling and changes if the photon properties change.
    \item The **absolute CMB temperature** is an observational constraint that remains fixed regardless of the theoretical scenario.
    \item In reality, the photon mass is constrained to be negligibly small by multiple observations, making this scenario purely pedagogical.
  \end{itemize}

\end{enumerate}


% ─────────────────────────────────────────────────────────────────────────────
